\sscp{Kayeli}{10-01-01}
Kayeli is a recently extinct language \citep{Grimes2010b} that
was spoken in some of the easternmost coastal communities of Buru.
\citet[20]{Collins1983} lists Kayeli as its own branch of Nunusaku (roughly our \tsc{Ambon-Seram}).
\citet[8]{Stresemann1927} describes Kayeli as a bridging language
between Sub-Buru and Sub-Ambon.
In our analysis  Kayeli groups with \tsc{Ambon-Seram} on the
basis of merging *ŋ/*n/*ñ > \it{n} (\srf{07-04-01-01},
\trf{07-08}), the merger of *d/*z/*R/*l/*j > \it{l} (\srf{10-01}),
and *p > \it{h}.

In \tsc{Sula-Buru} *k = **k,
but there is general weakening of *k found across \tsc{Ambon-Seram} languages.
Such weakening occurs in Kayeli,
along with lexically specific exceptions in which *k = \it{k}.
It is shown in example \qf{10-01},
alongside Buru cognates in which *k = \it{k} in all cases.
There was no clear contrast heard in Kayeli between vowel onsets
and glottal onsets from the last three speakers of Kayeli who
provided data to Grimes.

\noindent{\wg\st{6pt}\begin{tabularx}{\linewidth}{>{\hspace{-\tabcolsep}}l<{\hspace{-\tabcolsep}}>{\hspace{-\tabcolsep}\enspace}lll>{\raggedright\arraybackslash}X}
\lsptoprule
%\tbx{00-1}	&	PMP/PCM	&	Kayeli	&	Buru	&	gloss	\\	\midrule
%\endfirsthead
\tbx{10-01}	&	PMP/PCM	&	Kayeli	&	Buru	&	gloss	\\	\midrule
%\endhead
	&	*\tbr{k}a\tbr{k}a	&	\tbr{k}a\tbr{k}a-n	&	\tbr{k}ai, \tbr{k}a\tbr{k}a	&	elder sibling (same sex)	\\
	&	**\tbr{k}eRa	&	\tbr{k}ela	&	\tbr{k}eha	&	go up, rise	\\
	&	*bu\tbr{k}u	&	bu\tbr{k}u, bu\tbr{k}u-n	&	fo\tbr{k}o-n	&	joint	\\
	&	*bi\tbr{k}u, *biŋ\tbr{k}u	&	embi\tbr{k}o	&	efni\tbr{k}un	&	bent, crooked (Buru /en-fiku-n/ `elbow')	\\
	&	*\tbr{k}uh\tbr{k}uh	&	\tbr{k}o\tbr{k}o-ni	&	\tbr{k}u\tbr{k}u	&	fingernail 	\\
	&	*ti\tbr{k}aR	&	atin	&	\tbr{k}atin	&	sleeping mat (consonant met.)	\\
	&	*\tbr{k}ahiw	&	au	&	\tbr{k}au	&	tree, wood (Kayeli loan from \tsc{Sula-Buru})	\\
	&	*\tbr{k}utu	&	oto	&	\tbr{k}oto	&	head louse	\\
	&	*i\tbr{k}uR	&	i\tbr{k}o-i	&	i\tbr{k}u-n	&	tail	\\
	&	**malu\tbr{k}u	&	emloo	&	emlo\tbr{k}o	&	eel (see also \citealp[84]{Collins1983})	\\
	&	*pani\tbr{k}i	&	enhiiʔ	&		&	bat 	\\
	&	*hi\tbr{k}an	&	ian	&	i\tbr{k}an	&	fish	\\
	&	*la\tbr{k}aw	&	lawa	&		&	go, walk 	\\
	&	**waŋ\tbr{k}a	&	waa	&	waga	&	boat, canoe	\\
\lspbottomrule\end{tabularx}\mbox{}\\}

In \tsc{Sula-Buru} *p = **p,
but in Kayeli *p > \it{h},
consistent with the general weakening of *p > \it{f,
h,} {\0} in \tsc{Ambon-Seram}.
This is shown in example \qf{10-02}.
%
%\noindent{\wg\st{6pt}\begin{tabularx}{\linewidth}{>{\hspace{-\tabcolsep}}l<{\hspace{-\tabcolsep}}>{\hspace{-\tabcolsep}\enspace}lll>{\raggedright\arraybackslash}X}
%\lsptoprule
%%\tbx{00-2}	&	PMP	&	Kayeli	&	Buru	&	gloss	\\	\midrule
%%\endfirsthead
%\tbx{10-02}	&	PMP	&	Kayeli	&	Buru	&	gloss	\\	\midrule
%%\endhead
	%&	*\tbr{p}uqun	&	au \tbr{h}oni	&	\tbr{p}uun	&	tree	\\
	%&	*\tbr{p}aniki	&	en\tbr{h}iiʔ	&	(tagrihit)	&	fruitbat 	\\
	%&	*\tbr{p}a\tbr{p}an	&	\tbr{h}a\tbr{h}an-e	&	\tbr{p}a\tbr{p}an	&	board, plank	\\
	%&	**\tbr{p}a\tbr{p}u	&	\tbr{h}a\tbr{h}u	&	\tbr{p}a\tbr{p}u	&	boil something	\\
	%&	*\tbr{p}ajay	&	\tbr{h}ala	&	\tbr{p}ala	&	rice	\\
	%&	*\tbr{p}aRa `shelf' 	&	\tbr{h}ala-n	&	\tbr{p}aha-n	&	nest 	\\
	%&	*\tbr{p}aRih	&	\tbr{h}ali	&	\tbr{p}ahi	&	stingray	\\
	%&	**\tbr{p}eba	&	\tbr{h}eba	&	\tbr{p}efa	&	burn, bake	\\
	%&	*\tbr{p}əñu	&	\tbr{h}eno	&	\tbr{p}eno	&	sea turtle	\\
	%&	*\tbr{p}usuq, *\tbr{p}usuŋ	&	\tbr{h}oso-n	&	\tbr{p}oso-n	&	heart (Buru: `lower chest')	\\
	%&	**\tbr{p}una	&	\tbr{h}una	&	\tbr{p}una	&	do, make	\\
	%&	*\tbr{p}usəj	&	\tbr{h}use-n	&	\tbr{p}use-n	&	navel, belly button	\\
	%&	*h\<in\>i\tbr{p}i	&	em-ni\tbr{h}i	&	em-ni\tbr{p}i	&	dream	\\
	%&	*ha\tbr{p}uy	&	a\tbr{h}u-t	&	(bana)	&	fire 	\\
	%&	*qa\tbr{p}uR	&	a\tbr{h}ul-e	&	a\tbr{p}u	&	lime, calcium	\\
	%&	*Ra\tbr{p}us	&	la\tbr{h}o	&	ha\tbr{p}u	&	tie, bind	\\
	%&	*ni\tbr{p}ay, *ani\tbr{p}a	&	ni\tbr{h}a	&	(inewet)	&	snake 	\\
	%&	**sa\tbr{p}a	&	sa\tbr{h}a-n	&	sa\tbr{p}a-n	&	what? 	\\
	%&	**te\tbr{p}u' (PAMS)	&	te\tbr{h}u	&	te\tbr{p}u-t	&	chicken	\\
%\lspbottomrule\end{tabularx}\mbox{}\\}

{\wg\st{6pt}\begin{xltabular}[l]{\linewidth}
{>{\hspace{-\tabcolsep}}l<{\hspace{-\tabcolsep}}
>{\hspace{-\tabcolsep}\enspace}lll
>{\raggedright\arraybackslash}X}
\lsptoprule
%\tbx{00-2}	&	PMP	&	Kayeli	&	Buru	&	gloss	\\	\midrule
%\endfirsthead
	&	PMP	&	Kayeli	&	Buru	&	gloss	\\	\midrule
\endhead
\tbx{10-02}		&	*\tbr{p}uqun	&	au \tbr{h}oni	&	\tbr{p}uun	&	tree	\\
	&	*\tbr{p}aniki	&	en\tbr{h}iiʔ	&	(tagrihit)	&	fruitbat 	\\
	&	*\tbr{p}a\tbr{p}an	&	\tbr{h}a\tbr{h}an-e	&	\tbr{p}a\tbr{p}an	&	board, plank	\\
	&	**\tbr{p}a\tbr{p}u	&	\tbr{h}a\tbr{h}u	&	\tbr{p}a\tbr{p}u	&	boil something	\\
	&	*\tbr{p}ajay	&	\tbr{h}ala	&	\tbr{p}ala	&	rice	\\
	&	*\tbr{p}aRa `shelf' 	&	\tbr{h}ala-n	&	\tbr{p}aha-n	&	nest 	\\
	&	*\tbr{p}aRih	&	\tbr{h}ali	&	\tbr{p}ahi	&	stingray	\\
	&	**\tbr{p}eba	&	\tbr{h}eba	&	\tbr{p}efa	&	burn, bake	\\
	&	*\tbr{p}əñu	&	\tbr{h}eno	&	\tbr{p}eno	&	sea turtle	\\
	&	*\tbr{p}usuq, *\tbr{p}usuŋ	&	\tbr{h}oso-n	&	\tbr{p}oso-n	&	heart (Buru: `lower chest')	\\
	&	**\tbr{p}una	&	\tbr{h}una	&	\tbr{p}una	&	do, make	\\
	&	*\tbr{p}usəj	&	\tbr{h}use-n	&	\tbr{p}use-n	&	navel, belly button	\\
	&	*h\<in\>i\tbr{p}i	&	em-ni\tbr{h}i	&	em-ni\tbr{p}i	&	dream	\\
	&	*ha\tbr{p}uy	&	a\tbr{h}u-t	&	(bana)	&	fire 	\\
	&	*qa\tbr{p}uR	&	a\tbr{h}ul-e	&	a\tbr{p}u	&	lime, calcium	\\
	&	*Ra\tbr{p}us	&	la\tbr{h}o	&	ha\tbr{p}u	&	tie, bind	\\
	&	*ni\tbr{p}ay, *ani\tbr{p}a	&	ni\tbr{h}a	&	(inewet)	&	snake 	\\
	&	**sa\tbr{p}a	&	sa\tbr{h}a-n	&	sa\tbr{p}a-n	&	what? 	\\
	&	**te\tbr{p}u' (PAMS)	&	te\tbr{h}u	&	te\tbr{p}u-t	&	chicken	\\
\lspbottomrule
\end{xltabular}\addtocounter{table}{-1}}

The devoicing in Kayeli of *-mp-,
*-mb-, *ma-b is shown in example \qf{10-03},
contra \citeauthor{Blust2009}'s (\citeyear[72]{Blust2009})
claims about the “postnasal voicing
of stops” in such clusters in this region.
Note that PMP *b = \it{b} in Kayeli,
so the devoicing of the cluster is noteworthy.

\noindent{\wg\st{6pt}\begin{tabularx}{\linewidth}{>{\hspace{-\tabcolsep}}l<{\hspace{-\tabcolsep}}>{\hspace{-\tabcolsep}\enspace}lll>{\raggedright\arraybackslash}X}
\lsptoprule
%\tbx{10-03}	&	PMP	&	Kayeli	&	Buru	&	gloss	\\	\midrule
%\endfirsthead
\tbx{10-03}	&	PMP	&	Kayeli	&	Buru	&	gloss	\\	\midrule
%\endhead
	&	**(\tbr{ma-})\tbr{p}utu{\footnotemark}	&	\tbr{mp}utu	&	\tbr{p}oto 	&	Buru `hot to touch'; Kayeli `fever'	\\
	&	*t-u\tbr{mp}u	&		&	to\tbr{b}o-n	&	master, lord	\\
	&	*u\tbr{mp}u	&	o\tbr{p}o-ni	&		&	grandparent, grandchild	\\
	&	*\tbr{ma-b}u(Rr)it	&	\tbr{p}ori-n	&	\tbr{b}ohi-n	&	buttocks	\\
	&	**la\tbr{mb}u	&	la\tbr{p}u-n	&	la\tbr{b}u-n	&	shirt 	\\
	&	*Ru\tbr{mb}ia	&	\tbr{mp}ia	&	\tbr{b}ia	&	sago palm	\\
	&	**so\tbr{mb}a	&	so\tbr{p}	&	so\tbr{b}a, so\tbr{b}a-k	&	sail (v); go by sea	\\
\lspbottomrule\end{tabularx}\mbox{}\\}
\footnotetext{For **ma-putu compare South Nuaulu \it{putu} `hot'
(blood) and Proto-Rote-Meto *mbutu `hot' \citep[273]{Edwards2021}.
Buru \it{poto} `hot to touch' is from **putu without initial
*ma-.}

Kayeli also groups with \tsc{Ambon-Seram} by inflecting its verbs
and having different inflectional paradigms,
some of which involve morphophonemic alternations.
For example, \it{k-ino} ‘\tsc{1sg}-drink’,
\it{n-ino} ‘\tsc{3sg}-drink’ (< PMP *inum ‘drink’; contrast with
Buru-Lisela \it{ino} for all person and number combinations).
It also shows morphophonemic alternations in \it{tese luma} ‘we (incl.)
stay home’, \it{irese luma} /i{\nbhyp}tese/ ‘\tsc{3sg}-stays home’ (versus Buru
\it{eptea huma} ‘sit home’ or \it{defo di huma} ‘stay \tsc{dist}
home’ = ‘stay at home’ invariant for all person and number combinations).

Apparent innovations in the lexicon of \it{au} ‘tree,
wood’, \it{wai-n} ‘younger sibling (same sex)’,
\it{waa} ‘boat’, \it{lama-ni} ‘eye’,
\it{haa} ‘four’, \it{nee} ‘six’ that follow \tsc{Sula-Buru} patterns
are attributable to contact with Lisela, Buru and Ambelau.

In the other direction,
Lisela has many loans from Kayeli as would be expected in a junior-senior
role in their political histories.
These include: \it{aba-t} ‘thatch’,
\it{ahul} ‘lime, chalk’, \it{buu} ‘defecate’, \it{em-see} ‘cold,
flu’, \it{foli} ‘conch shell, trumpet’,
\it{hala} ‘rice’, \it{hali} ‘stingray’,
\it{heno} ‘sea turtle’, \it{hosi} ‘citrus’,
\it{kewel} ‘deaf’, \it{sii} ‘enter’,
\it{tiso} ‘newborn’, and \it{ubaa} ‘crocodile’.


